\section[Categorías fibradas y categorías prefibradas]{Categorías fibradas y
categorías prefibradas. Productos y cambio de base en ellas}
\label{VI.6}

\begin{definition}
\label{VI.6.1}
Una categoría $\cal{F}$ sobre~$\cal{E}$ se llama una
\emph{categoría fibrada}
\index{fibrada (categoría)|hyperpage}%
(y se dice entonces que el funtor $\cal{F}\to\cal{E}$ es un
\emph{fibrado})
\index{fibrado (funtor)|hyperpage}%
si satisface los dos axiomas siguientes:
\begin{itemize}
\item{$\Fib_{\mathrm{I}}$}
\marginpar{165}
Para todo morfismo $f\colon T\to S$ en~$\cal{E}$, existe el funtor de
imagen inversa por~$f$ en~$\cal{F}$.
\item{$\Fib_{\mathrm{II}}$} El compuesto de dos morfismos cartesianos es
cartesiano.
\end{itemize}
Una categoría $\cal{F}$ sobre~$\cal{E}$ que satisface la condición
$\Fib_{\mathrm{I}}$ se llama una \emph{categoría prefibrada sobre~$\cal{E}$}.
\index{prefibrada (categoría)|hyperpage}%
\end{definition}

Si~$\cal{F}$ es una categoría fibrada (\resp prefibrada) sobre~$\cal{E}$,
una subcategoría $\cal{G}$ de~$\cal{F}$ se llama una
\emph{subcategoría fibrada}
\index{subcategoría fibrada (\resp prefibrada)|hyperpage}%
(\resp una \emph{subcategoría prefibrada}) si es una categoría fibrada
(\resp prefibrada) sobre~$\cal{E}$ y, además, el funtor de inclusión es
cartesiano. Si, por ejemplo, $\cal{G}$ es una subcategoría \emph{plena}
de~$\cal{F}$, se ve que esto significa que, para todo morfismo
$f\colon T\to S$ en~$\cal{E}$ y todo $\xi\in\Ob(\cal{G}_S)$,
$f^*_{\cal{F}}(\xi)$ es $T$-isomorfo a un objeto de~$\cal{G}_T$. Otro
caso interesante es el siguiente: siendo $\cal{F}$ una categoría fibrada
sobre~$\cal{E}$, consideremos la subcategoría $\cal{G}$ de~$\cal{F}$ que
tiene los mismos objetos y cuyos morfismos son los morfismos
\emph{cartesianos} de~$\cal{F}$; en particular, los morfismos de
$\cal{G}_S$ son los isomorfismos de~$\cal{F}_S$. Se ve inmediatamente que
es, en efecto, una subcategoría fibrada de~$\cal{F}$, pues en la biyección
$$
\Hom_T(\eta',\eta) \isomto \Hom_f(\eta',\xi)
$$
relativa a un $f$-morfismo cartesiano $\alpha$ en~$\cal{F}$, a los
$T$-isomorfismos del primer miembro corresponden los morfismos cartesianos
del segundo. Por definición, las secciones cartesianas
$\cal{E}\to\cal{F}$
\ifthenelse{\boolean{orig}}{(corresponden}{corresponden}
entonces biunívocamente a los $\cal{E}$-funtores cualesquiera
$\cal{E}\to\cal{G}$ (pero obsérvese que el funtor natural
$$
\SheafHom_{\cal{E}/-}(\cal{E},\cal{G})\to
\SheafHom_\cart(\cal{E},\cal{F})=
\varprojLim(\cal{F}/\cal{E})
$$
es fiel, pero en general no es plenamente fiel, es decir, no es un
isomorfismo).

\begin{remarquesstar}
Sea $\cal{F}$ una categoría sobre~$\cal{E}$. Las condiciones siguientes
son equivalentes:~(i)~Todos los morfismos de~$\cal{F}$ son cartesianos;
(ii)~$\cal{F}$ es una categoría fibrada sobre~$\cal{E}$ y las
$\cal{F}_S$ son grupoides%
\ifthenelse{\boolean{orig}}{,}{}
(es decir, todo morfismo en~$\cal{F}_S$ es un isomorfismo). Se dice
entonces que $\cal{F}$ es una categoría \emph{fibrada en grupoides}
\index{fibrada en grupoides|hyperpage}%
\marginpar{166}
sobre~$\cal{E}$. Son estas las que se encuentran sobre todo en la
\og teoría de los espacios de módulos\fg. Si $\cal{E}$ es un grupoide, se
demuestra que las condiciones~(i) y~(ii) equivalen también a la siguiente:
(iii)~$\cal{F}$ es un grupoide y el funtor de proyección
$p\colon\cal{F}\to\cal{E}$ es transportable (\cf \Ref{VI.4.4}). Por
ejemplo, si $\cal{E}$ y $\cal{F}$ son grupoides tales que
\ifthenelse{\boolean{orig}}
{$\Ob \cal{E}$ y $\Ob \cal{F}$}
{$\Ob(\cal{E})$ y $\Ob(\cal{F})$}
se reducen a un punto, de modo que $\cal{E}$ y $\cal{F}$ quedan definidos,
salvo isomorfismo, por grupos $E$ y $F$, y el funtor
$p\colon\cal{F}\to\cal{E}$ queda definido por un homomorfismo de grupos
$p\colon F\to E$, entonces $\cal{F}$ está
\ifthenelse{\boolean{orig}}{fibrado}{fibrada}
sobre~$\cal{E}$ si y solo si $p$ es sobreyectivo, es decir, si $p$ define
una \emph{extensión} del grupo $E$ por el grupo $G=\Ker p$.
\end{remarquesstar}

\begin{proposition}
\label{VI.6.2}
Sea $F\colon\cal{F}\to\cal{G}$ una $\cal{E}$-equivalencia. Para que
$\cal{F}$ sea una categoría fibrada (\resp prefibrada) sobre~$\cal{E}$,
es necesario y suficiente que $\cal{G}$ lo sea.
\end{proposition}

Se deduce fácilmente de las definiciones y de la observación señalada más
arriba, según la cual un morfismo
\ifthenelse{\boolean{orig}}{$\alpha x$}{$\alpha$}
en~$\cal{F}$ es cartesiano si y solo si $F(\alpha)$ lo es.

\begin{proposition}
\label{VI.6.3}
Sean $\cal{F}_1$, $\cal{F}_2$ dos categorías sobre~$\cal{E}$, y sea
$\alpha=(\alpha_1,\alpha_2)$ un morfismo en
$\cal{F}=\cal{F}_1\times_{\cal{E}}\cal{F}_2$. Para que $\alpha$ sea
cartesiano, es necesario y suficiente que sus componentes lo sean.
\end{proposition}

Sea, en efecto, $\xi_i$ el destino y $\eta_i$ la fuente de~$\alpha_i$, y
sea $f\colon T\to S$ el morfismo de~$\cal{E}$ tal que $\alpha_1$ y
$\alpha_2$ sean $f$-morfismos. Para todo
$\eta'=(\eta'_1,\eta'_2)$ en~$\cal{F}_T$, se tiene un diagrama conmutativo
$$
\xymatrix{ \Hom_T(\eta',\eta) \ar[r] \ar[d] & \Hom_f(\eta',\xi) \ar[d]
\\ \Hom_T(\eta'_1,\eta_1)\times\Hom_T(\eta'_2,\eta_2) \ar[r] &
\Hom_f(\eta'_1,\xi_1)\times \Hom_f(\eta'_2,\xi_2) }
$$
donde las flechas verticales son biyecciones. Por tanto, si una de las
flechas horizontales es una biyección, también lo es la otra. Esto muestra
ya que, si $\alpha_1$, $\alpha_2$ son cartesianos (y, por tanto, la segunda
flecha horizontal es biyectiva), entonces $\alpha$ lo es. El recíproco se
ve tomando en el diagrama anterior $\eta'_i=\eta_i$, de donde
$\Hom_T(\eta'_i,\eta_i)\neq\emptyset$: primero para $i=2$, lo que prueba
que $\alpha_1$ es cartesiano, y luego para $i=1$, lo que prueba que
$\alpha_2$ es cartesiano.

\begin{corollaire}
\label{VI.6.4}
Sea
\marginpar{167}
$\cal{F}=\cal{F}_1\times_{\cal{E}}\cal{F}_2$, y sea $F=(F_1,F_2)$ un
$\cal{E}$-funtor $\cal{G}\to\cal{F}$. Para que
\ifthenelse{\boolean{orig}}{$\cal{F}$}{$F$}
sea cartesiano, es necesario y suficiente que $F_1$ y $F_2$ lo sean. Se
obtiene así un isomorfismo de categorías
$$
\SheafHom_\cart(\cal{G}, \cal{F}_1 \times_{\cal{E}} \cal{F}_2) \isomto
\SheafHom_\cart(\cal{G},\cal{F}_1)
\times\SheafHom_\cart(\cal{G},\cal{F}_2)
$$
y, en particular (tomando $\cal{G}=\cal{E}$), un isomorfismo de categorías
$$
\varprojLim(\cal{F}_1 \times_{\cal{E}}
\cal{F}_2/\cal{E}) \isomto
\varprojLim(\cal{F}_1/\cal{E}) \times
\varprojLim(\cal{F}_2/\cal{E})
$$
\end{corollaire}

\begin{corollaire}
\label{VI.6.5}
Sean $\cal{F}_1$ y $\cal{F}_2$ dos categorías fibradas (\resp prefibradas)
sobre~$\cal{E}$; entonces su producto fibrado
$\cal{F}=\cal{F}_1\times_{\cal{E}}\cal{F}_2$ es una categoría fibrada
(\resp prefibrada) sobre~$\cal{E}$.
\end{corollaire}

Estos resultados se extienden, por lo demás, al producto fibrado de una
familia cualquiera de categorías sobre~$\cal{E}$.

\begin{proposition}
\label{VI.6.6}
Sean $\cal{F}$ una categoría sobre~$\cal{E}$, con funtor de proyección
$p$, y $\lambda\colon\cal{E}'\to\cal{E}$ un funtor; consideremos
$\cal{F}'=\cal{F}\times_{\cal{E}}\cal{E}'$ como una categoría
sobre~$\cal{E}'$ mediante el funtor de proyección
$p'=p\times_{\cal{E}}\id_{\cal{E}'}$. Sea $\alpha'$ un morfismo
de~$\cal{F}'$. Para que $\alpha'$ sea un morfismo cartesiano, es necesario
y suficiente que su imagen $\alpha$ en~$\cal{F}$ lo sea.
\end{proposition}

La demostración es inmediata y se deja al lector.

\begin{corollaire}
\label{VI.6.7}
Para todo funtor cartesiano $F\colon\cal{F}\to\cal{G}$ de categorías
sobre~$\cal{E}$, el funtor $F'=F\times_{\cal{E}}\cal{E}'$ de
$\cal{F}'=\cal{F}\times_{\cal{E}}\cal{E}'$ en
$\cal{G}'=\cal{G}\times_{\cal{E}}\cal{E}'$ es cartesiano.
\end{corollaire}

Por consiguiente, el funtor $\SheafHom_{\cal{E}}(\cal{F},\cal{G})\to
\SheafHom_{\cal{E}'}(\cal{F}',\cal{G}')$ considerado en
\ifthenelse{\boolean{orig}}{\ignorespaces}{el}
\No \Ref{VI.3} induce un funtor
$$
\SheafHom_\cart(\cal{F},\cal{G})\to
\SheafHom_\cart(\cal{F}',\cal{G}') \rm{;}
$$
en otras palabras, para $\cal{F}$, $\cal{G}$ fijas, \emph{puede
considerarse}
$$
\SheafHom_\cart(\cal{F}\times_{\cal{E}}\cal{E}',
\cal{G}\times_{\cal{E}}\cal{E}')
$$
\emph{como
\marginpar{168}
un funtor en~$\cal{E}'$, de la categoría
${\Cat_{/\cal{E}}}^\circ$ en~$\Cat$}. Si se hacen variar asimismo
$\cal{F}$, $\cal{G}$, se obtiene%
\ifthenelse{\boolean{orig}}{,}{}
un funtor de la categoría ${\Cat_{/\cal{E}}}^\circ\times\big(
\Cat^\cart_{/\cal{E}}\big)^\circ\times\Cat^\cart_{/\cal{E}}$
en~$\Cat$. Al tener en cuenta el isomorfismo
$$
\Hom_{\cal{E}'}(\cal{F}',\cal{G}') \isomto
\Hom_\cal{E}(\cal{F}\times_\cal{E}\cal{E}',\cal{G})
$$
considerado en el~\No \Ref{VI.3}, los $\cal{E}'$-funtores cartesianos
de~$\cal{F}'$ en~$\cal{G}'$ corresponden a los $\cal{E}$-funtores
$\cal{F}\times_\cal{E}\cal{E}'\to\cal{G}$ que transforman todo morfismo
cuya primera proyección es un morfismo cartesiano de~$\cal{F}$ en un
morfismo cartesiano de~$\cal{G}$. Tomando $\cal{F}=\cal{E}$, se obtiene
(tras un cambio de notación):

\begin{corollaire}
\label{VI.6.8}
$\varprojLim(\cal{F}'/\cal{E}')$ es isomorfa a la subcategoría plena de
$\SheafHom_{\cal{E}/-}(\cal{E}',\cal{F})$ formada por los
$\cal{E}$-funtores $\cal{E}'\to\cal{F}$ que transforman los morfismos
cualesquiera en morfismos cartesianos.
En particular, si $\cal{F}$ es una categoría fibrada y
$\tilde{\cal{F}}$ es la subcategoría de~$\cal{F}$ cuyos morfismos son
los morfismos cartesianos de~$\cal{F}$, entonces se tiene una biyección
$$
\Ob \varprojLim(\cal{F}'/\cal{E}') \isomto
\Hom_{\cal{E}/-}(\cal{E}',\tilde{\cal{F}}) \rm{.}
$$
\end{corollaire}

Esto precisa la manera en que la expresión
$\varprojLim(\cal{F}\times_{\cal{E}}\cal{E}'/\cal{E}')$ debe considerarse
como un funtor en~$\cal{E}'$ y en~$\cal{F}$, de la categoría
${\Cat_{/\cal{E}}}^\circ\times\Cat^\cart_{/\cal{E}}$ en la categoría
$\Cat$. Más adelante se verá una dependencia funtorial más completa con
respecto a~$\cal{E}'$, cuando se exija que $\cal{F}$ sea una categoría
fibrada.

\begin{corollaire}
\label{VI.6.9}
Sea $\cal{F}$ una categoría fibrada (\resp prefibrada) sobre~$\cal{E}$;
entonces $\cal{F}'=\cal{F}\times_{\cal{E}}\cal{E}'$ es una categoría
fibrada (\resp prefibrada) sobre~$\cal{E}'$.
\end{corollaire}

\begin{proposition}
\label{VI.6.10} Sean $\cal{F}$, $\cal{G}$ categorías prefibradas
sobre~$\cal{E}$, y sea $F$ un $\cal{E}$-funtor cartesiano de~$\cal{F}$
en~$\cal{G}$. Para que $F$ sea fiel,
\marginpar{169}
\ifthenelse{\boolean{orig}}
{\resp plenamente fiel, (\resp una $\cal{E}$-equivalencia)}
{(\resp plenamente fiel, \resp una $\cal{E}$-equivalencia)}
es necesario y suficiente que, para todo
\ifthenelse{\boolean{orig}}
{$S\in\Ob \cal{E}$,}
{$S\in\Ob(\cal{E})$,}
el funtor inducido $F_S\colon\cal{F}_S\to\cal{G}_S$ sea fiel
(\resp plenamente fiel, \resp una equivalencia).
\end{proposition}

Demostración inmediata a partir de las definiciones.

Para terminar este número, damos algunas propiedades de las categorías
fibradas que utilizan el axioma $\Fib_{\mathrm{II}}$.
\begin{proposition}
\label{VI.6.11} Sea $\cal{F}$ una categoría prefibrada sobre~$\cal{E}$.
Para que $\cal{F}$ sea fibrada, es necesario y suficiente que satisfaga la
condición siguiente:

$\Fib_{\mathrm{II}}'$: Sea $\alpha\colon\eta\to\xi$ un morfismo
cartesiano en~$\cal{F}$ sobre el morfismo $f\colon T\to S$ de~$\cal{E}$.
Para todo morfismo $g\colon U\to T$ en~$\cal{E}$ y todo
\ifthenelse{\boolean{orig}}
{$\zeta\in\Ob \cal{F}_U$,}
{$\zeta\in\Ob(\cal{F}_U)$,}
la aplicación $u\mto\alpha\circ u$:
$$
\Hom_g(\zeta,\eta)\to\Hom_{fg}(\zeta,\xi)
$$
es biyectiva.
\end{proposition}

En otras palabras, en una categoría \emph{fibrada} sobre~$\cal{E}$, los
diagramas cartesianos se caracterizan por una propiedad a priori más
fuerte que la de la
\ifthenelse{\boolean{orig}}
{definición, (que se obtiene tomando $g=\id_T$ en}
{definición (que se obtiene tomando $g=\id_T$ en}
el enunciado precedente).
\begin{corollaire}
\label{VI.6.12} Sean $\cal{F}$ una categoría sobre~$\cal{E}$ y $\alpha$
un morfismo en~$\cal{F}$. Para que $\alpha$ sea un isomorfismo, es
necesario que $p(\alpha)=f$ sea un isomorfismo y que $\alpha$ sea
cartesiano; el recíproco es cierto si $\cal{F}$ está fibrada
sobre~$\cal{E}$.
\end{corollaire}

En efecto, si $\alpha$ es un isomorfismo, evidentemente también lo es
$f=p(\alpha)$; para todo
\ifthenelse{\boolean{orig}}
{$\eta'\in\Ob\cal{F}_T$,}
{$\eta'\in\Ob(\cal{F}_T)$,}
la aplicación $u\mto\alpha\circ u$
$$
\Hom(\eta',\eta)\to\Hom(\eta',\xi)
$$
es biyectiva; como $f$ es un isomorfismo, se ve inmediatamente
\ifthenelse{\boolean{orig}}{que un}{que un}
elemento del primer miembro es un $T$-morfismo si y solo si su imagen en el
segundo es un $f$-morfismo; se obtiene, por tanto, una biyección
$$
\Hom_T(\eta',\eta)\to\Hom_f(\eta',\xi)
$$
lo
\marginpar{170}
que prueba la primera afirmación. Recíprocamente, supongamos que $f$ es un
isomorfismo y que $\alpha$ satisface la condición enunciada en
\ifthenelse{\boolean{orig}}{$\Fib_{\mathrm{II}'}$}{$\Fib_{\mathrm{II}}'$}
(lo que significa, cuando $\cal{F}$ está
\ifthenelse{\boolean{orig}}
{fibrado}
{fibrada}
sobre~$\cal{E}$, que $\alpha$ es cartesiano); entonces se ve
inmediatamente que, para todo
\ifthenelse{\boolean{orig}}
{$\zeta\in\Ob\cal{F}$}
{$\zeta\in\Ob(\cal{F})$}%
, la aplicación $u\mto\alpha\circ u$ de~$\Hom(\zeta,\eta)$ en
$\Hom(\zeta,\xi)$ es biyectiva; por tanto, $\alpha$ es un isomorfismo.
\begin{corollaire}
\label{VI.6.13} Sean $\alpha\colon\eta\to\xi$ y
$\beta\colon\zeta\to\eta$ dos morfismos componibles en la categoría
$\cal{F}$ fibrada sobre~$\cal{E}$. Si $\alpha$ es cartesiano, entonces
$\beta$ lo es si y solo si $\alpha\beta$ lo es.
\end{corollaire}

Se utiliza la definición de los morfismos cartesianos en la forma reforzada
de~\Ref{VI.6.11}.
